一个\textbf{算术模型}就是算术语言~$\Lang{L_A}$ 的一个结构。有一个特殊的此类模型，即\textbf{标准模型}~$\Struct{N}$，其论域 $\Domain{N}
= \Nat$，并且对 $\Obj 0$、$\prime$、$+$、$\times$ 和 $<$ 的解释分别由 $0$、后继函数、$\Nat$ 上的加法和乘法函数以及小于关系给出。$\Struct{N}$ 是理论 $\Th{Q}$ 和~$\Th{PA}$ 的模型。

更一般地，$\Lang{L_A}$ 的一个结构称为\textbf{标准的}，当且仅当它同构于~$\Struct{N}$。两个结构称为同构的，如果它们之间存在一个\textbf{同构}，即一个保持常项符号、函数符号和谓词符号解释的双射函数。根据\textbf{同构定理}，同构的结构是\textbf{初等等价的}，亦即它们满足相同的语句。在标准模型中，论域恰好就是所有数字~$\num{n}$ 的值构成的集合。

与~$\Struct{N}$ 不同构的 $\Th{Q}$ 和 $\Th{PA}$ 的模型称为\textbf{非标准的}。在非标准模型中，论域并不被数字的值所穷尽。一个元素 $x \in \Domain{M}$，若对于所有~$n \in \Nat$ 都有 $x \neq \Value{\num{n}}{M}$，则称为~$\Struct{M}$ 的一个\textbf{非标准元素}。如果 $\Sat{M}{\Th{Q}}$，那么非标准元素也必须遵守~$\Th{Q}$ 的公理，例如，它们有唯一的后继，可以相加、相乘，并可以用~$<$ 进行比较。$\Struct{M}$ 中的所有标准元素都（按 $\Assign{<}{M}$）小于所有非标准元素。非标准模型的存在性由紧致性定理保证，并且对于 $\Th{Q}$ 而言，可以相对容易地显式给出这样的模型。这类模型可用来证明，例如，$\Th{Q}$ 不足以证明某些语句，例如 $\Th{Q} \Proves/ \lforall[x][\lforall[y][\eq[(x+y)][(y+x)]]]$。其方法是定义一个非标准的 $\Struct{M}$，在其中非标准元素不遵守交换律。

$\Struct{PA}$ 的非标准模型则无法如此轻易地显式指定。通过证明 $\Th{PA}$ 能证明某些特定的语句，我们可以研究 $\Th{PA}$ 的一个非标准模型的非标准部分的结构。如果 $\Th{PA}$ 的一个非标准模型~$\Struct{M}$ 是可数的，那么每个非标准元素都属于某个由非标准元素组成的``块''，块中元素在~$\Assign{<}{M}$ 下的序类似于~$\Int$。这些块本身的排列则类似于~$\Rat$，也就是说，没有最小或最大的块，并且在任意两个块之间总存在另一个块。

任何可数模型都同构于一个论域为~$\Nat$ 的模型。如果在这样一个模型中，$\prime$、$+$、$\times$ 和 $<$ 的解释是可计算函数（关系），我们就称它为一个\textbf{可计算模型}。标准模型~$\Struct{N}$ 是可计算的，因为 $\Nat$ 上的后继、加法、乘法函数和小于关系都是可计算的。可以定义 $\Th{Q}$ 的可计算非标准模型，但 $\Struct{N}$ 是 $\Th{PA}$ 唯一的可计算模型。这就是\textbf{Tennenbaum 定理}。